% Converted from: manuscript/draft_v5_en/front/10_abstract.md
% Abstract --- body only; main.tex wraps this in \begin{abstract} ... \end{abstract}.

The functional safety engineering frameworks of the automotive industry --- headed by the ISO 26262 V-Model --- assume that the behaviour of a component is derived from a specification written in advance, is verified against expected outputs and is validated statically before deployment. Components trained by reinforcement learning (RL) break all three assumptions: their behaviour emerges from a stochastic optimisation, their output does not admit the classical notion of a \enquote{correct answer}, and their robustness outside the training distribution can only be characterised empirically.

This thesis proposes and evaluates an adapted V-Model built on five explicit modifications (A1--A5) that reintegrate an RL component into a traceable development cycle, consistent with ISO 26262, ISO 21448 (SOTIF), ISO/IEC TR 5469 and UL 4600, without giving up the bidirectional specification$\leftrightarrow$V\&V correspondence that gives the standard its value. The framework is instantiated on a lane-following case study with a 1:14 scale vehicle (CobraFlex) trained in the Gazebo simulator: an end-to-end front-camera policy --- a convolutional network that learns perception as well as control --- contained by a safety cage of six deterministic rules that operates on its own vision-based lane estimator, independent of the network. A second instantiation on a privileged state vector acts as a control arm, in order to isolate the effect of the cage and to quantify the cost of perception.

The central contribution is methodological: an executable life cycle whose bidirectional traceability is verified by software along the chain Hazard $\rightarrow$ Safety Requirement $\rightarrow$ Cage Rule $\rightarrow$ Scenario $\rightarrow$ Metric $\rightarrow$ Logged Evidence $\rightarrow$ Verdict, so that no gate review can be closed with orphan identifiers on either side.

The reference validation campaign executes 27 scenarios $\times$ 2 modes = 1,890 runs on the \code{complex\_b} circuit, in \enf{} mode (cage active) and \mon{} mode (cage observes but does not act), with the two-dimensional action policy (steering and throttle) selected by driving quality rather than by reward. The result supports three claims. First: inside the operational domain not a single contact with the road edge is recorded with the cage active, against 60 that the policy commits on its own when the cage only observes; outside the domain the number is reduced to less than half of the previous instantiation. Second: the cage is latent in its safety rules inside the domain --- it does not trigger \cagerule{01}/\cagerule{02}/\cagerule{03}/\cagerule{05} in clean nominal driving --- and it becomes the operative mechanism exactly where perception degrades or the vehicle leaves the domain. Third, and less comfortable: the rate limiter \cagerule{06}, formally a smoothness rule of a non-critical class, turns out to be the rule that keeps the vehicle in the lane on the tightest curves, to the point that without the cage the same policy leaves the road in 17 out of 25 endurance runs; the dependency is measured, and its origin --- co-adaptation to the limiter inside the training loop --- is declared as inferred. That transfer risk was declared before any hardware was touched, and the first measurement on the real vehicle does not confirm it: by selecting the checkpoint for independence from the cage rather than for reward, the limiter intervenes in 3.4\,\% of the cycles on hardware against 3.0\,\% in simulation. This is a single run and therefore does not settle the question.

The work also reports a negative verdict: the rule composition requirement (\sr{010}) is not satisfied --- under simultaneous co-activation of the lateral and heading rules, violations of the lateral margin persist --- and training a better policy halves the finding without changing its nature, which places it as a design problem of the cage and not of the policy. It is declared as a limitation and as future work instead of being reconciled. Finally, the sim-to-real gap is characterised in steps of increasing fidelity\sidenote{Gazebo $\rightarrow$ Isaac Sim $\rightarrow$ physical platform.}, with the physical step executed as a bring-up and closed without a results campaign. That bring-up already produces five things that simulation could not produce, four of them independent of the campaign: an inherited and falsified assumption --- the effective field of view of the camera is \qdeg{77.89} and not the \qdeg{90} that the simulator replicated, an error that scales every lateral quantity on which the cage acts and that no additional amount of simulation would have revealed, because the simulator had inherited it; the negative result that the policy validated in simulation does not transfer, because it had memorised the 6.5:1 handedness of the training circuit as a steering bias; the preliminary observation\sidenote{One run.} that a retraining oriented towards transfer covers \qm{18.05} of the real circuit in a single segment without triggering any safety rule; the localisation of a failure mode that simulation has no way of exhibiting --- the lane estimator fails depending on the place on the circuit and not on the action of the vehicle, measured with true-position metrology; and a class of defect that only appears with a vehicle in front of you, in which rules that are correct by specification have no defined operational behaviour and sensors whose failure mode enters directly into the cage's only speed input. No scenario has been scored on hardware and the cage has never modified an action on the real vehicle, so the physical evidence is bring-up evidence and not verdict evidence --- the physical campaign is declared not executed and remains as future work, and the gap table keeps its physical column as provisional --- and it is declared as such.

The result is an end-to-end traceable case study that provides evidence of the feasibility, the adoption cost and the limits of the proposed framework.

\bigskip

\noindent Keywords: reinforcement learning; autonomous driving; lane following; safety cage; functional safety; ISO 26262; SOTIF; SE4AI; scenario-based validation; traceability; sim-to-real gap; PPO; ROS2; Gazebo; CobraFlex.
